% LSPV-4 REPORT::RESPSEL1::2026-10-03T01:46:00-07:00
\documentclass[11pt]{article}
\usepackage[margin=.85in]{geometry}
\usepackage{amsmath,amssymb}
\begin{document}
\section*{LSPV-4 RESPSEL1: total NP-response selector closure}
\textbf{STATUS: OPEN.} P versus NP remains open. This report proves a counterfactual closure theorem and uses it as a hostile red-team test. It does not prove a SAT lower bound.

\section*{CORPUS REVIEWED}
The run surveyed the current repository surfaces and read root \texttt{main.tex}, \texttt{AGENTS.md}, \texttt{CODEX.md}, the repository pnp-research skill, TASK-001 through TASK-005, verified MAIN 0058, current LSPV-3 predecessor CERTREG1, current LSPV-4 report surface, governance, and the active HEC/CRS-PPC/CCSPC and product-span branches surfaced by the living synthesis. Historical LSPV-4 controls ACQ1, REREAD1, PCLOSE1-R1, RANDCLOSE1, REACQ1, PADINV1, ISOCAN1, WEIGHTPAD1, TRACECLOSE1, and LEGALFLAT1 were reconciled through the current corpus state. No missing failed scientific body was reconstructed.

The immediate predecessor LSPV-3 report CERTREG1 is verified: it has the required header, substantive body, \verb|\end{document}|, and final sentinel. Its monotonic certified-exponent guard is accounting only and does not alter the present causal red-team target.

\section*{Target and exact quantifiers}
The target is CRS-B / causal witness production. Let
\[
R(x,h,c,y)\in\{0,1\}
\]
be a deterministic polynomial-time predicate. Here \(x\) is the fixed instance, \(h\) is the polynomial-length prior causal history, \(c\) is the current challenge, and \(y\) is the stage response. Assume a polynomial \(p\) such that every legal tuple \((x,h,c)\) has at least one response \(y\) of length at most \(p(|x|)\):
\[
\forall (x,h,c)\in {\cal L}\;\exists y,\quad |y|\le p(|x|),\quad R(x,h,c,y)=1.
\]
The legality domain \({\cal L}\) may itself depend on the prior history. The theorem concerns only legal tuples.

\section*{RESPSEL1 theorem}
\textbf{Theorem.} If \(P=NP\), then there is a deterministic polynomial-time selector \(S\) such that for every legal tuple
\[
R(x,h,c,S(x,h,c))=1.
\]
Moreover, if there are at most \(T(|x|)\) adaptive stages for a polynomial \(T\), and the next challenge is supplied externally after the previous response, the stage responses can be generated online in total polynomial time.

\paragraph{Proof.}
For a legal tuple \(z=(x,h,c)\), define the prefix-extension language
\[
L_R=\{(z,u):\exists v,\ |uv|\le p(|x|),\ R(z,uv)=1\}.
\]
Because \(R\) is polynomial-time decidable and the completion \(v\) is polynomially bounded, \(L_R\in NP\). Under \(P=NP\), \(L_R\in P\). Starting from the empty prefix, query whether a valid completion exists with next bit \(0\); if yes choose \(0\), otherwise choose \(1\). Include a polynomially encoded length delimiter, or equivalently test all allowed lengths, to handle variable response length. After at most \(p(|x|)+O(\log p(|x|))\) prefix decisions, this produces a valid \(y\). Thus \(S\in FP\).

For polynomially many adaptive stages, after receiving \(c_i\), compute
\[
y_i=S(x,h_{i-1},c_i)
\]
and append \((c_i,y_i)\) to the history. The history length remains polynomial and each selector call is polynomial, so the entire online response process is polynomial. \(\square\)

\section*{Hostile consequence for HEC / CRS-PPC / CCSPC}
RESPSEL1 closes a specific proposed escape: unpredictability of future challenges plus NP-verifiable, polynomially balanced, total stage responses cannot by itself force a superpolynomial response cost under the \(P=NP\) counterfactual. The responder need not predict future challenges. It waits for the current challenge and computes a valid response by the selector above.

Therefore a surviving CRS-B theorem must require more than
\[
\text{fresh challenge}+\text{polynomial witness}+\text{polynomial verification}.
\]
It needs an independently proved SAME\_REALIZATION resource condition not reducible to merely finding some valid NP response. In particular, requiring the exact reference history, exact provenance, exact prime tag, or exact internal reconstruction is not legitimate unless ordinary static SAT is first compiled to that same required realization with all acquisition and recomputation costs charged.

\section*{Relation to TRACECLOSE1 and LEGALFLAT1}
This is not TRACECLOSE1. TRACECLOSE1 says polynomial-time functionals of a deterministic polynomial execution transcript remain polynomial. RESPSEL1 does not assume the desired response is already a polynomial-time trace functional; under the counterfactual \(P=NP\), it derives an FP selector from total polynomially balanced NP search.

This is also stronger in a different direction than LEGALFLAT1. LEGALFLAT1 flattens a polynomial-round legal game when a polynomial-size polynomial-time evaluator for the strategy is already supplied. RESPSEL1 shows that for total NP-verifiable stage-response relations, the \(P=NP\) counterfactual itself supplies a polynomial-time response selector. It does not assert that one selector wins arbitrary adversarial games whose winning condition couples future rounds globally.

\section*{SU/vMS/Hadamard and reconstruction audit}
RESPSEL1 does not alter the fixed-realization algebraic questions. Exact source typing, Booleanization, characteristic-two product-span/common-kernel growth, sparse-\(Q\) normalization, source-matched reconstruction parameters, and the PAIRDEG1/FULLPAIR6 asymptotic condition remain OPEN or CONDITIONAL.

The theorem does constrain acquisition claims. If a proposed reconstruction object is merely a polynomially balanced witness to a polynomial-time relation and existence is guaranteed at each legal stage, then under \(P=NP\) it is polynomially selectable. To survive, acquisition must force a resource not captured by such an NP search relation, or the contradiction must arise from a separately proved model-preserving compiler.

CERTREG1, ENVELOPE1, RAMDOM1, and CROSSDOM1 remain accounting results. RESPSEL1 supplies no numerical reconstruction lower bound.

\section*{PRIME-LQP / provenance audit}
PRIME-LQP cannot obtain hardness solely by defining a polynomial-size provenance response that is efficiently checkable and guaranteed to exist. Under \(P=NP\), RESPSEL1 selects some valid response in polynomial time. A demand for the unique historical response or SAME\_REALIZATION response is load-bearing and must be justified by an arbitrary-algorithm transfer theorem rather than by provenance terminology.

Polynomial prime/CRT annotations, hashes, deterministic transcript tags, and source labels remain compatible with \(P=NP\) whenever the relevant valid-response relation satisfies the theorem's hypotheses.

\section*{Wrappers and easy-family controls}
DEAD and COPY can alter history without changing SAT semantics. DELAY can enlarge a trace while preserving polynomiality. INTERP can change implementation while preserving decision behavior. RECOMP can regenerate a valid response rather than retain the reference response. RESPSEL1 is therefore implementation-insensitive at the level of existence of some valid response.

Horn, dual-Horn, 2-SAT, affine/Tseitin, and bounded-treewidth families reinforce the same guard: a causal protocol must not charge an exact preferred history when an easy solver can produce another valid response cheaply. Independent products, padding, renaming, and auxiliary-variable encodings likewise block representation-dependent charging unless invariance is proved.

\section*{Barrier audit}
RESPSEL1 relativizes in the following sense: if \(P^O=NP^O\) and \(R\) is polynomial-time decidable relative to the same oracle \(O\), the prefix-extension language is in \(NP^O=P^O\), so the selector is in \(FP^O\). Hence RESPSEL1 cannot be the nonrelativizing hinge.

The proof is black-box search-to-decision and provides no nonalgebrizing hinge. Current algebrization-barrier results therefore remain a separate guard. Natural Proofs is not directly triggered because RESPSEL1 defines no large constructive circuit property.

\section*{Primary-literature preflight}
Baker--Gill--Solovay's relativization barrier remains the governing oracle test. Ren--Williams, ECCC TR26-118 (2026), prove near-maximum circuit lower bounds for an exponential-time class with promise-MA queries and analyze bounded adaptive NP-query rounds and witness lengths; this is a model-specific circuit-lower-bound framework, not an unrestricted static-SAT-to-causal compiler. Li--Ren--Zhong, ECCC TR26-196 (2026), obtain many proof-complexity generators inside one demi-bits generator and restricted range-avoidance consequences; these do not force arbitrary SAT deciders to expose a preferred provenance response. Garlik--Gryaznov--Ren--Tzameret, ECCC TR26-133 (2026), prove strong Weak Rank Principle lower bounds in PCR over \(\mathbb F_2\), but the proof-system model remains load-bearing. Chen--Hu--Ren's algebrization-barrier work remains a warning that nonrelativization alone would not suffice.

\section*{DUPLICATION / CONTRADICTION AUDIT}
The corpus already contains the causal NDTM guard that a nondeterministic machine may wait for each challenge and guess a witness, and LEGALFLAT1 already closes supplied succinct strategies. RESPSEL1 is the deterministic counterfactual sharpening: assuming \(P=NP\), total polynomially balanced NP stage search has an FP selector, even when challenges arrive adaptively. This does not contradict the existing causal core; it narrows the missing CRS-B theorem.

No contradiction with CERTREG1 or the structural product-span program was found. No claim that fixed proof-system, communication, streaming, pebbling, or rank lower bounds transfer to unrestricted SAT is adopted.

\section*{WHAT CHANGED RELATIVE TO ROOT MAIN}
Root MAIN states that stagewise nondeterminism is an escape and that LEGALFLAT1 is a local closure. RESPSEL1 adds a sharper counterfactual deterministic closure: under \(P=NP\), every total polynomially balanced P-verifiable stage-response relation admits an FP selector, and polynomially many causal stages remain polynomial. This should be added as a hostile control, not as progress toward a separation.

\section*{ADOPTED / REJECTED / RETIRED PRIOR COMPONENTS}
Adopted at declared scope: TRACECLOSE1, REACQ1, LEGALFLAT1, HEC/CRS-PPC/CCSPC restricted-model lower bounds, FULLPAIR6/PAIRDEG1 finite structural certificates, and ENVELOPE1-family accounting.

Rejected as an unrestricted-runtime hinge: fresh challenge unpredictability plus ordinary NP witness production; exact-reference-history charging without a compiler; proof-system-specific hardness without arbitrary-algorithm transfer; representation size without acquisition.

FALSIFIED/RETIRED as a generic route: any claim that polynomially many adaptive total NP-verifiable response stages are intrinsically superpolynomial under the \(P=NP\) counterfactual.

\section*{Components consistent with \(P=NP\)}
TRACECLOSE1, REACQ1, LEGALFLAT1, RESPSEL1, polynomially selectable provenance responses, fixed-realization SU/vMS algebra, and all currently verified reconstruction-accounting identities are consistent with \(P=NP\). No load-bearing contradiction follows from them.

\section*{One next falsification target}
Take the strongest current CRS-B response relation and write its exact predicate
\[
R(x,h,c,y)
\]
with its size bound and totality domain. If it is polynomial-time verifiable, polynomially balanced, and total on legal histories, RESPSEL1 closes witness-production hardness under \(P=NP\). If one hypothesis fails, identify that exact failure and test whether it is genuinely forced by ordinary static SAT rather than inserted by the causal protocol.

\section*{Strict status}
RESPSEL1 is \textbf{PROVED/EXACT} as a conditional closure theorem with the explicit hypothesis \(P=NP\). Its application to any particular CRS-B construction is \textbf{CONDITIONAL} until the exact response relation is transcribed. The static-to-causal compiler, arbitrary-algorithm acquisition, source-matched no-sharing, SU/vMS asymptotic structural growth, and nonrelativizing/nonalgebrizing hinge remain \textbf{OPEN}. P versus NP remains open.

\section*{References}
T. Baker, J. Gill, and R. Solovay, Relativizations of the P=?NP Question, SIAM Journal on Computing 4(4), 1975.

H. Ren and R. Williams, Near-Maximum Circuit Lower Bounds for Exponential Time with Merlin-Arthur Queries, ECCC TR26-118, 2026.

X. Li, H. Ren, and Y. Zhong, Many Proof Complexity Generators Inside One Demi-Bits Generator, ECCC TR26-196, 2026.

M. Garlik, S. Gryaznov, H. Ren, and I. Tzameret, The Weak Rank Principle: Lower Bounds and Applications, ECCC TR26-133, 2026.

L. Chen, Y. Hu, and H. Ren, New Algebrization Barriers to Circuit Lower Bounds via Communication Complexity of Missing-String, 2025.
\end{document}
% END_OF_VERIFIED_REPORT::LSPV-4::RESPSEL1
